\scanpage{195}
% archive sheet number: 194
% source line 1
\uline{Catégories abéliennes}

% source lines 2--15
Soit $\unsure{\underline{C}|\underline{\mathcal{C}}}$ une catégorie abélienne où tous les objets sont de longueur finie. Alors $\operatorname{Ind}\underline{C}$ [catégorie des foncteurs contravariants exacts à gauche de $\underline{C}$ dans \deleted{\unsure{Ens}} $(\mathrm{Ab})$] est une \uline{catégorie abélienne} à \uline{limites inductives exactes}, à générateurs de longueur finie [nous dirons : \uline{loc. de longueur finie}] dont $\underline{C}$ est la sous-catégorie des objets de longueur finie. Il y a ainsi \unsure{essentiellement} correspondance biunivoque entre catégories abéliennes à objets de longueur finie, et catégories abéliennes à limites inductives exactes à générateurs de longueur finie.

% source lines 16--21; insertion above line 16
Soit $\underline{C}$ une catégorie abélienne\inserted{à limites inductives exactes}, $X\in\underline{C}$, le \uline{socle} de $X$ est la \deleted{\illegible{1}} somme des sous-objets simples de $X$, soit $F_0(X)$. On désigne par $F_1(X)$ l'image inverse du socle de $X/X_0$ etc, et on \unsure{trouve} ainsi une filtration croissante. \deleted{\unsure{Si} $\underline{C}$ \unsure{est}}
